% GENERATED FILE. Edit canonical lessons, then run npm run generate:books.
% canonical-source-hash: fnv1a64:b914fa954e1684e1
% canonical-lessons: FR-C260-je-pense-que, FR-C260-donc, FR-C260-revision-opinion

\chapter{What I Think, and So — je pense que, donc}
\label{ch:fr-what-i-think}
\hlchaptermodality{260}

\begin{chapteropening}
\textbf{By the end of this chapter:} \emph{I can give an opinion with je pense que and à mon avis, and point a reason forward with donc or alors.}

Complete FR-C260-revision-opinion: an opinion with its reason, said both ways round, with parce que and with donc.
\end{chapteropening}

\section[je pense que, à mon avis]{je pense que — giving a view}

\label{lesson:FR-C260-je-pense-que}



\begin{quote}
Say \textquotedblleft{}to think\textquotedblright{}. (\emph{Penser.}) Then \textquotedblleft{}I think\textquotedblright{}. (\emph{Je pense.})

A story says what happened. Giving reasons starts somewhere else: with what you think about it.
\end{quote}

\subsection*{You'll want to know: je pense que}
\textbf{je pense que} — \textquotedblleft{}I think that\textquotedblright{}. It is the everyday opinion opener.

\begin{quote}
\textbf{Je pense que} le train est cher. — I think the train is expensive. \textbf{Je pense que oui.} — I think so. · \textbf{Je pense que non.} — I don't think so.
\end{quote}

The \emph{que} cannot be dropped, as English drops its \emph{that}. After it, the verb stays plain: \emph{je pense que c'est cher}. Only a denial, \emph{je ne pense pas que…}, makes careful French bend the verb that follows into a form a later chapter builds.

Two near neighbours do the same job. \textbf{Je crois que}, from \emph{croire}, sounds a shade less sure. \textbf{Je trouve que}, from \emph{trouver}, is \textquotedblleft{}I find that\textquotedblright{}, for a judgement you have made yourself: \emph{je trouve que le film est ennuyeux}.

To hand the opinion back to a friend: \emph{\textbf{Et toi, qu'est-ce que tu en penses ?}} — \textquotedblleft{}And you, what do you think of it?\textquotedblright{} The little \emph{en} means \textquotedblleft{}of it\textquotedblright{}; keep the whole question as one piece for now.

\begin{grammarlens}[title={Grammar Lens: à mon avis}]
You have \textbf{l'avis}, \textquotedblleft{}a notice, an opinion\textquotedblright{}. \textbf{À mon avis} is \textquotedblleft{}in my opinion\textquotedblright{}, and it puts your standpoint first:

\begin{quote}
\textbf{À mon avis}, le train est cher, mais le petit musée est sympathique.
\end{quote}

Relaxed speech says \textbf{pour moi} for the same thing. \emph{Je pense que} puts the opinion first; \emph{à mon avis} says whose it is before you give it.
\end{grammarlens}

\begin{cousinweb}
\textbf{penser} is inherited, through Old French \emph{penser}, from Latin \emph{pēnsāre}, \textquotedblleft{}to weigh carefully\textquotedblright{} (EtymDB 2.1), so a thought is a thing weighed. \textbf{avis} goes back to Old French \emph{avis}, built on \emph{vis}, \textquotedblleft{}a view\textquotedblright{}, which is inherited from Latin \emph{vīsus}, \textquotedblleft{}a look, an appearance\textquotedblright{} (EtymDB 2.1). \emph{À mon avis} is \textquotedblleft{}at my view\textquotedblright{}: how the matter looks from where I stand.
\end{cousinweb}

\subsection*{Your turn}
Say these aloud:

\begin{itemize}
\raggedright
  \item \emph{Je pense que le train est cher.}
  \item both short answers — \emph{Je pense que oui.} then \emph{Je pense que non.}
  \item hand it back — \emph{Et toi, qu'est-ce que tu en penses ?}
  \item \emph{À mon avis, le film est ennuyeux.}
\end{itemize}

\begin{tcolorbox}[breakable,colback=teal!4,colframe=teal!35!black,title={Before you move on}]
Say \textquotedblleft{}I don't think so\textquotedblright{}. (\emph{\textbf{Je pense que non.}}) Can the \emph{que} be dropped? (\textbf{No.}) Which opener says whose view it is first? (\emph{\textbf{À mon avis.}}) And what did Latin \emph{pēnsāre} mean? (\textbf{To weigh.})
\end{tcolorbox}
\section[donc, alors]{donc, alors — the arrow that points forward}

\label{lesson:FR-C260-donc}



\begin{quote}
Give an opinion with a reason: \textquotedblleft{}I think it's going to rain, because there's no sun.\textquotedblright{} (\emph{Je pense qu'il va pleuvoir, parce qu'il n'y a pas de soleil.})

\emph{Parce que} runs from the claim back to its reason. This lesson runs the other way.
\end{quote}

\begin{grammarlens}[title={Grammar Lens: the two arrows}]
The same two facts, in opposite orders:

\begin{quote}
J'ai pris le métro \textbf{parce qu'}il pleuvait. — I took the metro \textbf{because} it was raining. Il pleuvait, \textbf{donc} j'ai pris le métro. — It was raining, \textbf{so} I took the metro.
\end{quote}

\textbf{Parce que} points \textbf{back}, from what happened to its cause. \textbf{Donc} points \textbf{forward}, from the cause to what followed, which is why the cause has to come first. Which one you reach for depends on which fact you want to say first: a choice about emphasis, not about grammar.

The tenses keep their jobs. \emph{Il pleuvait} is the scene, the reason; \emph{j'ai pris} is the event, the consequence.
\end{grammarlens}

\begin{grammarlens}[title={Grammar Lens: alors, the chatty one}]
\textbf{alors} can stand in the same place: \emph{Il pleuvait, \textbf{alors} j'ai pris le métro.} It is lighter and more conversational, the \textquotedblleft{}so\textquotedblright{} of a friend telling you about their day. \textbf{Donc} is the neutral one, as at home in a letter as in speech.

\begin{quote}
chatty: \emph{alors} · neutral: \emph{donc}
\end{quote}

Both point forward. \emph{Alors} also keeps its first meaning, \textquotedblleft{}then, at that time\textquotedblright{}, so in a story it can do both jobs at once. And when \emph{donc} opens a clause, its \emph{c} is usually said: \emph{Donc, je reste.}
\end{grammarlens}

\begin{cousinweb}
\textbf{alors} is \emph{à} plus Old French \emph{lors}, \textquotedblleft{}then\textquotedblright{}, which goes back to Latin \emph{illā hōrā}, \textquotedblleft{}at that hour\textquotedblright{} (EtymDB 2.1): a word about time that became a word about consequence, as English \textquotedblleft{}then\textquotedblright{} did. \textbf{donc} is usually traced to Late Latin \emph{dunc}, \textquotedblleft{}then, so\textquotedblright{}, whose own making is debated; EtymDB 2.1 follows it only as far as Middle French \emph{donc}.
\end{cousinweb}

\subsection*{Your turn}
Say these aloud:

\begin{itemize}
\raggedright
  \item back — \emph{J'ai pris le métro parce qu'il pleuvait.}
  \item forward — \emph{Il pleuvait, donc j'ai pris le métro.}
  \item the same, to a friend — \emph{Il pleuvait, alors j'ai pris le métro.}
  \item \emph{Il faisait froid, alors je suis rentré.}
  \item \emph{À mon avis, le train est cher, donc je prends le métro.}
\end{itemize}

\begin{tcolorbox}[breakable,colback=teal!4,colframe=teal!35!black,title={Before you move on}]
Which way does \emph{parce que} point, and which way \emph{donc}? (\emph{Parce que} \textbf{back} to the cause; \emph{donc} \textbf{forward} to the consequence.) Which of the two forward words is the chatty one? (\emph{\textbf{Alors.}}) What did Latin \emph{illā hōrā} mean? (\textbf{At that hour.})
\end{tcolorbox}
\section[(review)]{Review — a view, and the arrows behind it}

\label{lesson:FR-C260-revision-opinion}



\begin{quote}
Say \textquotedblleft{}I think so\textquotedblright{}. (\emph{Je pense que oui.}) Then \textquotedblleft{}in my opinion\textquotedblright{}. (\emph{À mon avis.})
\end{quote}

\begin{grammarlens}[title={What you've built}]
Two lessons, and the start of every argument you will make in French. An opinion, with \emph{je pense que}, \emph{je crois que} or \emph{à mon avis}. A reason pointing back, with \emph{parce que}. A consequence pointing forward, with \emph{donc} or \emph{alors}.

Here is one view, built both ways:

\begin{quote}
\emph{Je prends le métro, parce que je pense que le train est cher.} \emph{Je pense que le train est cher, donc je prends le métro.}
\end{quote}

Same facts, same view. The first starts from the decision; the second starts from the opinion that caused it.
\end{grammarlens}

\subsection*{Your turn}
Say these aloud:

\begin{itemize}
\raggedright
  \item \emph{Je pense que le train est cher.}
  \item the reason, back — \emph{Je prends le métro parce que le train est cher.}
  \item the same, forward — \emph{Le train est cher, donc je prends le métro.}
  \item in the past — \emph{Il pleuvait, alors j'ai pris le métro.}
  \item a view of your own, with \emph{à mon avis} and one reason
\end{itemize}

\begin{tcolorbox}[breakable,colback=teal!4,colframe=teal!35!black,title={Before you move on}]
Which word points back to a cause? (\emph{\textbf{Parce que.}}) Which two point forward? (\emph{\textbf{Donc}} and \emph{\textbf{alors}}.) In \emph{il pleuvait, alors j'ai pris le métro}, which verb is the scene? (\emph{\textbf{Il pleuvait.}})
\end{tcolorbox}
